% amsthm must be loaded before the class: sn-jnl defines its theorem
% styles (thmstyleone/two/three) inside an \@ifpackageloaded{amsthm}
% guard evaluated at class-load time.
\RequirePackage{amsthm}
\documentclass[sn-mathphys,NameDate]{sn-jnl}

%% ============================================================
%% PAPER C (JoS candidate) — drafted 2026-06-11.
%%
%% Source of record: paper/poa/notes.md (PoA-1..5, hardened
%% 2026-06-11 pm) and paper/stochastic/notes.md (S1, S2). Every
%% theorem stated here is proved in those notes and cross-checked
%% against the exhaustive sweep (PriceOfAnarchy.scala, 39 configs);
%% the two bracketing results are machine-checked in Lean 4
%% (lean/AssemblyFlowShop/MakespanLowerBound.lean and
%% PriceOfAnarchy.lean, zero sorry).
%%
%% Numbering map (paper <-> notes):
%%   Thm 1  <-> Fact 0 + S1(a)            (exact optimum)
%%   Cor 2  <-> S1(c)                     (stochastic, fixed order)
%%   Thm 3  <-> S2                        (stochastic completeness)
%%   Thm 4  <-> PoA-1                     (full anarchy, makespan)
%%   Thm 5  <-> PoA-2                     (work-conserving, makespan)
%%   Thm 6  <-> PoA-4                     (full anarchy, flowtime)
%%   Thm 7  <-> PoA-5                     (work-conserving, flowtime)
%%   Prop 8 <-> PoA-3                     (WCT unbounded)
%% ============================================================

\usepackage{graphicx}
\usepackage{amsmath,amssymb,amsfonts}
\usepackage{amsthm}
\usepackage{mathrsfs}
\usepackage{xcolor}
\usepackage{textcomp}
\usepackage{manyfoot}
\usepackage{booktabs}

\theoremstyle{thmstyleone}
\newtheorem{theorem}{Theorem}
\newtheorem{lemma}[theorem]{Lemma}
\newtheorem{corollary}[theorem]{Corollary}
\newtheorem{proposition}[theorem]{Proposition}
\theoremstyle{thmstylethree}
\newtheorem{definition}[theorem]{Definition}
\theoremstyle{thmstyletwo}
\newtheorem{remark}[theorem]{Remark}

\raggedbottom

\begin{document}

\title[The Price of Local Ordering]{The Price of Local Ordering in
Two-Stage Assembly Flow Shops}

\author*[1]{\fnm{Eric} \sur{Bowman}}\email{eric.bowman@king.com}

\affil[1]{\orgname{King}, \orgaddress{\city{Berlin}, \country{Germany}}}

\abstract{In the divisible-work two-stage assembly flow shop,
permutation schedules---one job order shared by every machine---%
dominate pointwise, so adopting a shared order is never worse. This
paper quantifies the converse: how much can abandoning the shared
order cost? In the fluid homogeneous model we first pin the optimum
exactly: every schedule satisfies $k\,C_{\max} \ge W + w_{\max}$,
and the synchronized schedule attains it under \emph{every} job
order; consequently, under arbitrarily distributed random work, a
fixed shared order---chosen with no foresight---stochastically
dominates every scheduling policy on makespan, and the permutation
class is stochastically complete for all regular objectives. Against
this exact baseline the price of local ordering is sharp: arbitrary
per-machine orders combined with a downstream stage that may hold
capacity idle for a preferred job cost exactly
$2W/(W+w_{\max}) \to 2$ on makespan, while \emph{any}
work-conserving downstream stage---regardless of its processing
order---caps the price at $2 - 1/m$, where $m$ is the number of
stage-1 machines; both constants are tight. The same dichotomy
governs total flowtime, with tight constants $3$ and $3 - 2/m$; for
weighted completion time no constant-factor ceiling exists. The
structural conclusion is that the expensive downstream behavior is
\emph{work-withholding}, not misordering: a stage that idles while
released work waits forfeits as much again as all upstream
misalignment combined. Every constant is verified by exhaustive
enumeration over the instance family of the companion paper, where
the worst cases match the theory exactly, and the two bracketing
results are formalized in Lean~4.}

\keywords{assembly flow shop, permutation schedule, divisible work,
price of anarchy, work conservation, stochastic scheduling}

\pacs[MSC Classification]{90B35}

\maketitle

%% ============================================================
\section{Introduction}
\label{sec:intro}

Dominance theorems are one-sided. For the two-stage assembly flow
shop with divisible work---$m$ parallel stage-1 machines of $k$
identical workers feeding a single assembly stage of $k$ workers---%
for every schedule there is a \emph{permutation schedule}, with one
job order shared by all stage-1 machines, that completes every job
at the assembly stage no later \citep{bowmanA}. A companion paper
maps where such dominance survives on deeper stage structures
\citep{bowmanB}. These results say that coordinating on a shared
order is \emph{never worse}. They are silent on the question any
adopter of the policy actually asks: how much worse is it \emph{not}
to coordinate?

This paper answers that question exactly, in the fluid
(real-divisible) homogeneous form of the model. Three groups of
results.

\emph{First, the baseline is the true optimum} (Section~\ref{sec:opt}).
Every schedule---any per-machine orders, any work divisions, any
downstream discipline, preemptive or not, idle or not---satisfies
\[
  k \, C_{\max} \;\ge\; W + w_{\max},
\]
where $W$ is the total work and $w_{\max}$ the largest job, and the
synchronized schedule attains this bound under \emph{every} shared
order $\pi$ (Theorem~\ref{thm:optimum}). Two consequences follow.
Under arbitrarily distributed random work---any joint law---a fixed
shared order, chosen before any work is known, is pathwise
makespan-optimal, hence stochastically dominant over every policy
including clairvoyant ones: for the date the whole portfolio of jobs
lands, foresight is worth nothing and the choice of shared order is
worth nothing; only \emph{sharedness} matters
(Corollary~\ref{cor:fixed-order}). And the pointwise dominance
theorem of \citet{bowmanA} couples, realization by realization, into
a stochastic completeness statement: no policy beats the permutation
class on any regular objective in any stochastic sense
(Theorem~\ref{thm:completeness}). All prices below are therefore
measured against the exact optimum, not against a possibly
suboptimal synchronized baseline.

\emph{Second, the price of local ordering on makespan is an exact
dichotomy} (Section~\ref{sec:makespan}). If every stage-1 machine
chooses its own order and the downstream stage may \emph{withhold}
work---idle while released work waits, as a stage does when it
insists on processing in a preferred order whose next job has not
yet arrived---the worst-case makespan ratio is exactly
$2W/(W+w_{\max})$, approaching~$2$ (Theorem~\ref{thm:poa-full}).
If instead the downstream stage is \emph{work-conserving}---never
idle while released work is unfinished, but otherwise free to
process in any order whatsoever---the price drops to $2 - 1/m$,
and this is tight even against a clairvoyant downstream
(Theorem~\ref{thm:poa-wc}). The proof of the upper bound never uses
the downstream \emph{order}; only non-idleness. The structural
conclusion, which we believe is the practically important one, is
that for makespan the expensive downstream sin is withholding, not
misordering: a downstream stage that holds capacity idle for its
preferred job forfeits as much again as all upstream misalignment
combined.

\emph{Third, the same dichotomy governs flowtime, and no ceiling
exists for weighted completion time}
(Sections~\ref{sec:flowtime}--\ref{sec:wct}). For equal works the
total-flowtime price is exactly $(3n+1)/(n+3) \to 3$ under full
anarchy and $3 - 2/m$ under a work-conserving downstream, both
tight. For weighted completion time the price is unbounded:
concentrating value on one small job and ordering it locally
multiplies the weighted objective by $\Theta(W/w_{\max})$.

Every constant above is exact at finite sizes in a way that can be
checked by brute force, and we have checked it: over the
39-configuration exhaustive enumeration family of \citet{bowmanA}
(roughly $4.0 \times 10^{6}$ schedules), with the unsynchronized
population split by downstream discipline, the worst cases agree
with the theory to the digit, including two structural predictions
--- the two downstream classes must coincide once $m \ge n$, and the
withholding class must be empty at $n = 2$
(Section~\ref{sec:verification}). The two bracketing results
(Theorem~\ref{thm:optimum}'s lower bound, and the universal upper
bound behind Theorem~\ref{thm:poa-full}) are additionally formalized
in Lean~4 with no unproven steps.

\subsection{Related work}

The question belongs to the price-of-anarchy tradition of
\citet{koutsoupias1999}: compare a coordinated optimum against the
worst outcome of uncoordinated choices. Our setting differs from
game-theoretic scheduling in that we do not model incentives or
equilibria; ``anarchy'' here is simply the absence of a shared
order, which is how the phenomenon presents in the motivating
applications (portfolios of cross-team initiatives; multi-stage
processing pipelines). The two-stage assembly flow shop originates
with \citet{lee1993} and \citet{potts1995}, who proved permutation
dominance and approximation results in the unit-capacity case; see
\citet{hariri1997} for exact methods and \citet{bowmanA} for the
divisible-work generalization and the literature between.

The constant $2 - 1/m$ of Theorem~\ref{thm:poa-wc} has the same
form as Graham's bound for list scheduling on $m$ identical
machines \citep{graham1966,graham1979}, and the resemblance is not
coincidental: both results bound the damage of arbitrary orderings
by a busy-period argument in which work conservation is the only
property used. The settings are different---Graham's $m$ counts
parallel machines executing the jobs, our $m$ counts upstream
machines feeding an assembly stage, and our ratio is measured
against an exactly known optimum---but the mechanism, that
non-idleness alone buys a factor-$(2-1/m)$ guarantee, appears to be
the same phenomenon in two guises. The stochastic results connect
to the classical observation that simple index policies retain
optimality under uncertainty \citep{rothkopf1966,smith1956,
pinedo2016}; Corollary~\ref{cor:fixed-order} is an extreme instance
in which the optimal policy requires no distributional information
at all. Divisible (malleable) work is surveyed by
\citet{drozdowski2009}.

%% ============================================================
\section{Model and policy classes}
\label{sec:model}

There are $n$ jobs; job $j$ requires $w_j > 0$ units of work on
each of $m$ stage-1 machines and $w_j$ units on a single stage-2
(assembly) machine. Every machine has $k$ identical unit-rate
workers. Work is divisible: a machine that devotes all $k$ workers
to one job processes it at rate $k$. Write
$W = \sum_j w_j$ and $w_{\max} = \max_j w_j$. The
\emph{assembly constraint}: job $j$ may not be processed at stage~2
before all $m$ stage-1 machines have completed it; the time this
happens is $j$'s \emph{availability} $a_j$. All jobs are present at
time zero. This is the homogeneous fluid form of the model of
\citet{bowmanA}; the integer-divisible form used there for the
pointwise theorem appears here only in
Theorem~\ref{thm:completeness} and in the computational section.

A schedule is \emph{team-sequential} if every machine processes one
job at a time at full rate $k$, stage-1 machines without inserted
idleness. (Stage-1 idleness and unbalanced work divisions can only
increase availabilities and are briefly discussed after each upper
bound; the lower bound of Theorem~\ref{thm:optimum} holds with no
restriction at all.) Within team-sequential schedules we distinguish
the downstream discipline:

\begin{definition}[Policy classes]\label{def:classes}
A team-sequential schedule is:
\begin{itemize}
\item \emph{synchronized} if all stage-1 machines and the stage-2
machine process the jobs in one common order $\pi$;
\item \emph{work-conserving downstream} if the stage-2 machine is
never idle while some released job is unfinished---its processing
order is otherwise arbitrary;
\item \emph{work-withholding downstream} if it processes the jobs
back-to-back in some fixed preferred order, each job starting as
soon as the stage is free \emph{and} the job is available---so the
stage idles exactly when its order's next job has not yet arrived,
even while other released work waits. (This is the behavior of a
stage that enforces a priority list internally; it is the only
inserted idleness we allow downstream, and the enumeration model of
Section~\ref{sec:verification} matches it.)
\end{itemize}
\emph{Upstream anarchy} means the stage-1 machines use arbitrary,
possibly different orders. \emph{Full anarchy} means upstream
anarchy with a work-withholding downstream.
\end{definition}

Objectives are functions of the stage-2 completion vector
$(C_j)_{j}$: makespan $C_{\max} = \max_j C_j$, total flowtime
$\sum_j C_j$, and weighted completion time $\sum_j v_j C_j$ for
non-negative values $v_j$.

%% ============================================================
\section{The exact optimum, deterministically and stochastically}
\label{sec:opt}

\begin{theorem}[The optimum, and its order-independence]
\label{thm:optimum}
\mbox{}
\begin{enumerate}
\item[(a)] Every schedule---team-sequential or not, with arbitrary
work divisions, inserted idleness, and downstream discipline,
preemptive or not---satisfies
$k\,C_{\max} \ge W + w_{\max}$.
\item[(b)] For every order $\pi$, the synchronized team-sequential
schedule with order $\pi$ has
$C_{\max} = (W + w_{\max})/k$.
\end{enumerate}
Hence the optimal makespan is exactly $(W+w_{\max})/k$, attained by
the shared order---\emph{any} shared order.
\end{theorem}

\begin{proof}
(a) Let $M$ be a job of maximal work and let $a = a_M$ be its
availability; fix one stage-1 machine that completes $M$ at time
$a$ (if several, any). By time $a$ that machine has processed at
most $a k$ units of work, of which $w_{\max}$ belong to $M$. A job
is eligible for stage~2 only when complete on \emph{all} machines,
in particular on this one, so the total work of jobs other than $M$
available by time $a$ is at most $a k - w_{\max}$, and stage~2 can
have completed at most that much by time $a$. At least
$W - (a k - w_{\max})$ units of stage-2 work remain at time $a$,
processed at rate at most $k$:
\[
C_{\max} \;\ge\; a + \frac{W - a k + w_{\max}}{k}
\;=\; \frac{W + w_{\max}}{k},
\]
independent of $a$. The argument uses only capacity counting; it is
indifferent to work divisions, idleness, preemption, and the number
of machines.

(b) Under the shared order $\pi$, stage~1 completes the job at
position $s$ at time $S_s/k$ where
$S_s = \sum_{t \le s} w_{\pi(t)}$. Writing $F_s$ for its stage-2
completion and $M_s = \max_{t\le s} w_{\pi(t)}$, induction on $s$
gives $F_s = (S_s + M_s)/k$: indeed
$F_s = \max(S_s, \, S_{s-1} + M_{s-1})/k + w_{\pi(s)}/k
= (S_{s-1} + \max(M_{s-1}, w_{\pi(s)}))/k + w_{\pi(s)}/k
= (S_s + M_s)/k$. At $s = n$ this is $(W + w_{\max})/k$ for every
$\pi$.
\end{proof}

\begin{remark}[Lean formalization]\label{rem:lean}
Part (a) is machine-checked in Lean~4 over the concrete
integer-divisible machine model of \citet{bowmanA}, with stage~2
abstracted by a cumulative-progress structure whose axioms
(progress capped by the work, zero before release, total rate at
most $k$) are satisfied by every physical execution; the formal
statement, $k\,T \ge W + w_M$ for \emph{every} job $M$, is
correspondingly stronger than (a). The formalization also revealed
that the choice of machine in the proof is immaterial: any machine
bounds the released work, because every completion lies below the
availability. The development has no unproven steps and uses only
the standard axioms.
\end{remark}

The optimum being attained by every shared order has a stochastic
consequence that we state in full generality.

\begin{corollary}[A fixed shared order is stochastically optimal
for makespan]\label{cor:fixed-order}
Let the works $(W_1,\dots,W_n)$ be random with an arbitrary joint
distribution, and fix any order $\pi$ in advance. Then, pathwise,
the synchronized schedule with order $\pi$ achieves
$C_{\max} = (W + W_{\max})/k \le C_{\max}(S)$ for every scheduling
policy $S$, including policies with full knowledge of the
realization. In particular the fixed shared order stochastically
dominates every policy on makespan.
\end{corollary}

\begin{proof}
Apply Theorem~\ref{thm:optimum} realization by realization: (b)
gives the value of the fixed-order schedule, (a) bounds every
competitor below by the same value.
\end{proof}

For objectives beyond makespan a single fixed order no longer
suffices, but the \emph{class} of shared-order schedules remains
unbeatable, again under arbitrary randomness. The following
statement is in the integer-divisible model of \citet{bowmanA},
whose pointwise theorem it couples; it is the formal sense in which
all prices in this paper are measured against an optimum that
stochastic work cannot dislodge.

\begin{theorem}[Stochastic completeness of the permutation class]
\label{thm:completeness}
Let the works have an arbitrary joint distribution and let $S$ be
any scheduling policy, possibly randomized and clairvoyant. There
is a works-measurable permutation-schedule policy $S^*$ with
$C_j(S^*) \le C_j(S)$ for every job $j$, almost surely.
Consequently, for every coordinatewise-monotone objective $f$,
$f(C(S^*)) \le f(C(S))$ almost surely, and no policy beats the
permutation class in expectation, in any quantile, or in the usual
stochastic order, for any regular objective.
\end{theorem}

\begin{proof}
The construction of \citet{bowmanA} maps each realization of the
works and of $S$'s schedule to a permutation schedule dominating it
pointwise; the permutation used is obtained by sorting the
realized availabilities, so the map is measurable (there are
finitely many orders, and completion times are measurable in the
works). Composing $S$ with this map yields $S^*$; the domination is
pointwise on the same probability space, i.e.\ a coupling, and all
claimed comparisons follow.
\end{proof}

\begin{remark}
Theorem~\ref{thm:completeness} is existential in the order, which
may depend on the realization. Corollary~\ref{cor:fixed-order} shows
that for makespan the dependence vanishes. Whether a fixed
\emph{value-ordered} permutation is optimal in expectation for
weighted completion time under exchangeable works---the natural
analogue of weighted-shortest-expected-processing-time results
\citep{rothkopf1966}---is open; see Section~\ref{sec:discussion}.
\end{remark}

%% ============================================================
\section{The price of local ordering: makespan}
\label{sec:makespan}

\begin{theorem}[Full anarchy doubles the schedule, exactly]
\label{thm:poa-full}
Over team-sequential schedules with upstream anarchy and a
work-withholding downstream,
\[
\sup_{\text{instances, schedules}}
\frac{C_{\max}}{(W+w_{\max})/k} \;=\;
\sup_{\text{instances}} \frac{2W}{W + w_{\max}} \;=\; 2,
\]
with the inner supremum attained for every instance with $m \ge 2$:
no team-sequential schedule exceeds $2W/k$, and for every instance
some unsynchronized schedule reaches it.
\end{theorem}

\begin{proof}
\emph{Upper bound.} Stage-1 machines are work-conserving and
process at rate $k$, so every machine finishes all its work at
$W/k$; hence $a_j \le W/k$ for every $j$. A downstream stage that
processes its jobs back-to-back in any order $\sigma$, waiting only
for availabilities, satisfies the max-plus recursion
$F_{\sigma(s)} = \max(F_{\sigma(s-1)}, a_{\sigma(s)}) +
w_{\sigma(s)}/k$, whence by induction
$C_{\max} \le \max_j a_j + W/k \le 2W/k$.

\emph{Lower bound.} Fix any job $J$. Let stage-1 machine~1 process
$J$ last (so $a_J = W/k$) and machine~2 use any different order
(so the schedule is unsynchronized). Let the downstream stage
process $J$ \emph{first}: it idles until $W/k$, then has all $W$
units to do with no further idleness, finishing at exactly $2W/k$.
Against the optimum $(W+w_{\max})/k$ of
Theorem~\ref{thm:optimum} the ratio is $2W/(W+w_{\max})$, which
tends to $2$ as $w_{\max}/W \to 0$.
\end{proof}

\begin{remark}
The upper bound is also machine-checked in Lean~4 at the
release-time interface (makespan at most the largest availability
plus the total downstream work, for every downstream order), along
with the availability bound and computed witnesses for the smallest
anti-aligned instance. With unbalanced work divisions the
\emph{lower} bound only improves for the adversary; the sweep of
Section~\ref{sec:verification} shows split badness stacking on top
of the ordering price exactly where divisions are not forced
balanced.
\end{remark}

The construction above needs the downstream stage to \emph{wait}
for $J$. The next theorem shows this is the only way to get
anywhere near a factor of two.

\begin{theorem}[Work conservation caps the price at $2 - 1/m$]
\label{thm:poa-wc}
Over team-sequential schedules with upstream anarchy and ANY
work-conserving downstream discipline,
\[
C_{\max} \;\le\; \Bigl(2 - \frac{1}{m}\Bigr)\frac{W}{k}
+ \frac{w_{\max}}{k},
\]
and the constant is tight: for equal works $w$ and $m \mid n$ there
are stage-1 orders under which \emph{every} downstream
discipline---work-conserving, withholding, or clairvoyant---has
$C_{\max} \ge (2 - 1/m)\,nw/k + w/k$. Hence the price of upstream
anarchy under a work-conserving downstream is exactly $2 - 1/m$
asymptotically.
\end{theorem}

\begin{proof}
\emph{Upper bound.} Let $t_1$ be the last instant stage~2 is idle
before $C_{\max}$ (if stage~2's work is complete at its last idle
instant, then $C_{\max}$ is at most the last availability, at most
$W/k$, and we are done). At $t_1$, work conservation means every
available job is already finished at stage~2, so the remaining
stage-2 work is at most $W - w(A(t_1))$, where $A(t)$ is the set of
jobs available by $t$ and $w(\cdot)$ sums their works; stage~2 is
busy at rate $k$ on $[t_1, C_{\max}]$, giving
$C_{\max} \le t_1 + (W - w(A(t_1)))/k$.

Since stage~2 idles at $t_1$ with some unfinished job not yet
available, some stage-1 machine is still busy; stage-1 machines
have no upstream and are work-conserving, so that machine has
processed exactly $t_1 k \le W$. Every still-busy machine $i$ has
processed exactly $t_1 k$ units, of which at most $w_{\max}$ sit in
its single in-progress job, so its set $F_i(t_1)$ of finished jobs
satisfies $w(F_i(t_1)) \ge t_1 k - w_{\max}$ (machines already done
satisfy this trivially). A job is available iff finished on all $m$
machines, so by inclusion--exclusion for the work measure,
\[
w(A(t_1)) = w\Bigl(\bigcap_i F_i(t_1)\Bigr)
\;\ge\; \sum_i w(F_i(t_1)) - (m-1)W
\;\ge\; m(t_1 k - w_{\max}) - (m-1)W .
\]
Set $T^* = (m-1)W/(mk) + w_{\max}/k$. If $t_1 \le T^*$ then
$C_{\max} \le t_1 + W/k \le (2-1/m)W/k + w_{\max}/k$. If
$t_1 \ge T^*$ then
$C_{\max} \le t_1(1-m) + (mW + m\,w_{\max})/k$, which is decreasing
in $t_1$ and equals $(2-1/m)W/k + w_{\max}/k$ at $t_1 = T^*$.

\emph{Lower bound (rotation construction).} Equal works $w$,
$m \mid n$. Let machine $i$ process the jobs in the cyclic rotation
of $(1,\dots,n)$ by $(i-1)n/m$ positions. The $m$ rotation offsets
are evenly spaced, so job $j$'s positions across the machines form
the set $\{(j \bmod (n/m)) + r\,n/m + 1 : r = 0,\dots,m-1\}$; its
\emph{worst} position is at least $(m-1)n/m + 1$. Hence no job is
available before $T_0 = ((m-1)n/m + 1)\,w/k$, and stage~2---however
scheduled---must process all $W = nw$ units after $T_0$:
$C_{\max} \ge T_0 + nw/k = (2 - 1/m)\,nw/k + w/k$. Against the
optimum $(n+1)w/k$ the ratio tends to $2 - 1/m$. (For $m \nmid n$,
rotating by $\lfloor n/m \rfloor$ degrades the constant by
$O(w/k)$.)
\end{proof}

\begin{remark}[The dichotomy]\label{rem:dichotomy}
The upper bound of Theorem~\ref{thm:poa-wc} never refers to the
downstream \emph{order}; non-idleness is the only property used. So
for makespan the downstream behaviors split into exactly two
classes: any work-conserving stage, even one that ignores
priorities entirely, is protected at $2 - 1/m$; only a
work-withholding stage---one that holds capacity idle for its
preferred job---reaches the full price of
Theorem~\ref{thm:poa-full}. With two upstream machines the numbers
are stark: upstream misalignment alone costs at most $50\%$;
downstream withholding doubles it. In operational terms, enforcing
a priority order \emph{inside} a downstream stage is precisely a
withholding discipline whenever the preferred job has not yet
arrived, and it forfeits as much again as all upstream misalignment
combined. The resemblance of $2 - 1/m$ to Graham's list-scheduling
bound \citep{graham1966} is discussed in
Section~\ref{sec:intro}; both are busy-period consequences of work
conservation alone.
\end{remark}

%% ============================================================
\section{The price of local ordering: flowtime}
\label{sec:flowtime}

For flowtime we take equal works $w$ (so $W = nw$) and a
\emph{non-preemptive} one-job-at-a-time downstream stage; the
spacing step below fails under preemptive interleaving, whose
constant we conjecture is the same
(Section~\ref{sec:discussion}). The synchronized flowtime is
$\sum_s (s+1)w/k = n(n+3)w/(2k)$ by
Theorem~\ref{thm:optimum}(b), and it is optimal for equal works:
the earliest stage-2 completion is at least $2w/k$ (the first job
needs $w/k$ at each stage), and the non-preemptive downstream
completes jobs at least $w/k$ apart, so every schedule has
$C_{(s)} \ge (s+1)w/k$.

\begin{theorem}[Full anarchy triples flowtime, exactly]
\label{thm:poa-flow}
Over team-sequential schedules with a non-preemptive downstream,
for equal works,
\[
\sup \frac{\sum_j C_j}{n(n+3)w/(2k)} \;=\; \frac{3n+1}{n+3}
\;\longrightarrow\; 3 ,
\]
the supremum over upstream anarchy with a withholding downstream,
attained for every $n$ when $m \ge 2$.
\end{theorem}

\begin{proof}
\emph{Upper bound.} Stage~2 processes one job at a time, each
taking $w/k$, so the sorted completions satisfy
$C_{(s+1)} \ge C_{(s)} + w/k$; with
$C_{(n)} = C_{\max} \le 2nw/k$ (Theorem~\ref{thm:poa-full}),
$C_{(s)} \le (n+s)w/k$ and
$\sum_j C_j \le (3n^2+n)\,w/(2k)$ for \emph{every} team-sequential
schedule. The ratio to $n(n+3)w/(2k)$ is $(3n+1)/(n+3)$.

\emph{Lower bound.} In the construction of
Theorem~\ref{thm:poa-full}'s lower bound, stage~2 idles until
$nw/k$ and then completes one job per $w/k$: completions are
exactly $(n+s)w/k$ for $s = 1,\dots,n$, meeting the upper bound
term by term.
\end{proof}

\begin{theorem}[Work conservation caps the flowtime price at
$3 - 2/m$]\label{thm:poa-flow-wc}
Under the hypotheses of Theorem~\ref{thm:poa-flow} but with any
work-conserving downstream, for equal works and $m \mid n$,
\[
\sum_j C_j \;\le\;
\Bigl(\tfrac{3}{2} - \tfrac{1}{m}\Bigr)\frac{n^2 w}{k} + O(n),
\]
and the rotation construction attains
$\sum_j C_j = (\tfrac{3}{2} - \tfrac{1}{m})\,n^2 w/k + O(n)$ under
every downstream discipline. The price is exactly $3 - 2/m$
asymptotically.
\end{theorem}

\begin{proof}
\emph{Upper bound.} By Theorem~\ref{thm:poa-wc},
$C_{(n)} \le ((2-1/m)n + 1)\,w/k$; spacing gives
$C_{(s)} \le ((1-1/m)n + s + 1)\,w/k$; summing yields the claim.

\emph{Lower bound.} In the rotation instance the $m$ jobs with
residue $c$ become available exactly at slot $(m-1)n/m + c + 1$, so
$m$ new jobs arrive per slot from $T_0$ on while stage~2 completes
one per slot: it never starves after $T_0$, and completions are
exactly $T_0 + s\,w/k$, summing to
$n T_0 + n(n+1)w/(2k) = (\tfrac32 - \tfrac1m) n^2 w/k + O(n)$.
Since no job is available before $T_0$, no downstream discipline
does better than these completions term by term.
\end{proof}

The pattern of Remark~\ref{rem:dichotomy} thus persists in the
flowtime currency---$3 - 2/m$ for any work-conserving downstream
versus $3$ for a withholding one---with the same interpretation:
withholding forfeits as much again as all upstream misalignment.

%% ============================================================
\section{No ceiling for weighted completion time}
\label{sec:wct}

\begin{proposition}[The weighted price is unbounded]
\label{prop:wct}
For every $B$ there is an instance and an unsynchronized
team-sequential schedule whose weighted completion time exceeds $B$
times the optimum.
\end{proposition}

\begin{proof}
Put value $v$ on a single job $J$ of work $w$ and negligible value
on the others. The synchronized schedule with $J$ first completes
it at $2w/k$, so the optimum is at most
$2vw/k + o(v)$. The construction of Theorem~\ref{thm:poa-full}
applied to $J$ completes it at $2W/k$, giving weighted objective at
least $2vW/k$. The ratio is $W/w + o(1)$, unbounded as the number
of jobs grows.
\end{proof}

Schedule damage is capped at a factor of two
(Theorem~\ref{thm:poa-full}); value damage has no cap, and it grows
exactly with the concentration of value. Combined with
Corollary~\ref{cor:fixed-order} this separates the two roles of the
shared order: \emph{any} shared order is optimal for the date all
work lands; a \emph{correct} (value-aware) shared order is what
protects when value lands.

%% ============================================================
\section{Computational verification}
\label{sec:verification}

The constants above are exact at finite sizes, so they can be
confronted with exhaustive enumeration. We use the
39-configuration instance family of \citet{bowmanA}
($n \le 5$ jobs, $m \in \{2,3\}$ stage-1 machines,
$k \in \{2,3\}$ workers, symmetric, asymmetric, value-weighted and
machine-heterogeneous works; about $4.0 \times 10^{6}$ schedules in
total), with every combined schedule classified as synchronized or
unsynchronized and, within the unsynchronized population, the
downstream order classified as availability-consistent
(work-conserving in the integer model) or re-ordered (withholding).
The integer model forces balanced divisions exactly when $w_j = k$;
in those configurations the fluid constants must be---and
are---attained exactly.

\begin{table}[h]
\caption{Worst-case makespan price (per cent over the optimum) of
the unsynchronized population, split by downstream discipline, on
the division-forced configurations, against the finite-$n$ theory
of Theorems~\ref{thm:poa-full} and~\ref{thm:poa-wc}.}
\label{tab:split}
\begin{tabular}{@{}lllll@{}}
\toprule
Configuration & \multicolumn{2}{l}{Work-conserving} &
\multicolumn{2}{l}{Withholding} \\
 & observed & theory & observed & theory \\
\midrule
$n=3$, $m=2$ & $+25.0\%$ & $5/4$ & $+50.0\%$ & $6/4$ \\
$n=4$, $m=2$ & $+40.0\%$ & $7/5$ & $+60.0\%$ & $8/5$ \\
$n=5$, $m=2$ & $+33.3\%$ & $8/6$ & $+66.7\%$ & $10/6$ \\
$n=3$, $m=3$ & $+50.0\%$ & coincide & $+50.0\%$ & $6/4$ \\
$n=4$, $m=3$ & $+40.0\%$ & $7/5$ & $+60.0\%$ & $8/5$ \\
$n=2$, any $m$ & $+33.3\%$ & $4/3$ & \multicolumn{2}{l}{class
empty (forced)} \\
\botrule
\end{tabular}
\end{table}

Beyond the exact agreement of every entry, two structural
predictions are reproduced: at $n = 3$, $m = 3$ the two downstream
classes coincide, as they must once $m \ge n$ (the rotation
construction then defeats every downstream discipline); and at
$n = 2$ the withholding class is empty, as it must be
(availabilities tie, so every downstream order is
availability-consistent). The flowtime worst cases similarly match
Theorems~\ref{thm:poa-flow} and~\ref{thm:poa-flow-wc} exactly
where divisions are forced ($+100.0\% = (3n+1)/(n+3)$ at $n=5$;
$+50.0\%$ for the work-conserving class, the rotation value).
Configurations with division freedom exceed the balanced-split
constants in both classes, quantifying how division badness stacks
on top of ordering badness.

Averages tell the operational story: across the full family, the
mean makespan price of an unsynchronized schedule is $39.6\%$
(per-configuration means $20.0$--$63.2\%$), the mean flowtime price
$60.6\%$, and the mean weighted price $69.2\%$ (worst $+475.5\%$,
under value skew $100{:}1{:}1$); on average across configurations,
more than $99.5\%$ of unsynchronized schedules are strictly later
than the optimum on every objective (no configuration falls below
$95.6\%$). Within the unsynchronized population, conditioning
on an availability-consistent downstream roughly halves the mean
price (e.g.\ from $+53.1\%$ to $+26.3\%$ at $n=5$, $m=2$)---work
conservation alone recovers about half the damage of upstream
misalignment, before any upstream coordination.

The two bracketing results are formalized in Lean~4 over the
concrete machine model of \citet{bowmanA}: the universal lower
bound $k\,C_{\max} \ge W + w_M$ for every job $M$
(Theorem~\ref{thm:optimum}(a)), with stage~2 abstracted so that any
physical execution satisfies the axioms; and the universal upper
bound at the release-time interface behind
Theorem~\ref{thm:poa-full}, together with computed witnesses for
the smallest anti-aligned instance ($2W$ versus $W + w_{\max}$).
Both developments contain no unproven steps. Code, enumeration
outputs, and the Lean development are in the archived artifact (see
Declarations).

%% ============================================================
\section{Discussion}
\label{sec:discussion}

Together with \citet{bowmanA,bowmanB}, the results complete a
two-sided account of the shared order in this model. The dominance
theorem says adopting it is free; this paper says declining it
costs $40\%$ on average in the exhaustive family, at most a factor
of two on the schedule (exactly), without limit on concentrated
value---and that the two halves of the price have different owners.
Upstream misalignment is worth $2 - 1/m$; everything beyond is the
downstream stage's choice to hold capacity idle for a preferred
job. For practice the asymmetry is the message: a downstream stage
that simply keeps pulling whatever has arrived is protected by
Theorem~\ref{thm:poa-wc} no matter how badly it orders; the
behavior to eliminate is waiting---which is what enforcing a
priority list inside a stage operationally means whenever the
list's next item has not arrived. The stochastic results add that
none of this depends on knowing the work: a fixed shared order is
exactly optimal for the portfolio date under any randomness
(Corollary~\ref{cor:fixed-order}), while protecting \emph{value}
requires the order to be right
(Proposition~\ref{prop:wct})---sharedness is free insurance,
correctness buys value.

Open problems. (i) The flowtime spacing arguments require a
non-preemptive downstream; we conjecture the constants $3$ and
$3 - 2/m$ survive preemption. (ii) The flowtime constants are
proved for equal works; the general-works analogues should follow
the same pattern. (iii) The expected price under \emph{random}
(rather than adversarial) per-machine orders is strictly positive
and empirically large ($39.6\%$ mean above); its exact asymptotics
are open. (iv) Whether the value-ordered fixed permutation is
expectation-optimal for weighted completion time under exchangeable
works (the WSEPT analogue) is open. (v) The online version---jobs
arriving over time, the shared order maintained by preemptive
priority---is the natural continuation and the model closest to the
motivating applications.

%% ============================================================
\section*{Declarations}

\textbf{Funding.}
This research did not receive any specific grant from funding
agencies in the public, commercial, or not-for-profit sectors.

\textbf{Competing interests.}
The author declares no competing interests.

\textbf{Author contributions.}
The sole author conceived the work, developed the proofs, and wrote
the manuscript.

\textbf{Data availability.}
The enumeration outputs supporting
Section~\ref{sec:verification} are available in the repository
below.

\textbf{Code availability.}
The simulation code, exhaustive-enumeration suite, and the Lean~4
formalization are available at
\url{https://github.com/ebowman/synchronization-paper-artifact}
and archived at
\url{https://doi.org/10.5281/zenodo.20075388}.

\textbf{Generative AI disclosure.}
AI-assisted tooling (Anthropic Claude) was used for literature
search, drafting support, simulation code, and proof checking; all
mathematical content was verified by the author, by exhaustive
computational enumeration, and---where stated---by the Lean~4 proof
assistant.

\bibliography{price}

\end{document}
